\begin{recipe}
[% 
    preparationtime = {\unit[35]{min}},
    portion = {\portion{3}},
    %source = {},
    calory = {230kcal (590kcal) | 16g (28g) Protein | 23g (96g) KH | 9g (11g) Fett},
    price = {1,35€ (1,95€)},
    created = {14.03.2026},
    rating = {5},
]
{Bolognese Variante 2}

%\graph{% pictures
%    big=pic/hauptgericht/11_vegane-bolognese
%}

\introduction{%
Noch eine weitere Bolo, die finde ich noch ein bisschen besser unt intensiver schmeckt als die im Rezept davor.
}

\ingredients{
    120 g & Karotten\index{Gemüse!Karotte}\\
    2 & kleine Zwiebeln\\
    90 g & Sellerieknolle\index{Gemüse!Sellerie}\\
    40 g & Pastinake\index{Gemüse!Pastinake}\\
    45 g & Lauch\index{Gemüse!Lauch}\\
    650 g & Rispen-Tomaten\index{Gemüse!Tomate}\\
    80 g & getr. Sojageschnetzeltes\index{Tofu \& Fleischalternativen!Soja}\\
    2 EL & Tomatenmark\\
    4 EL & Sojasahne\index{Milchprodukte \& Alternativen!Sojasahne}\\
    4 EL & Olivenöl\\
    1 EL & Butter\\
    1 TL & Zucker\\
    \nicefrac{1}{2} TL & MSG\\ 
    1 TL & Oregano\\
    \nicefrac{1}{2} TL & Thymian\\
    opt. 1 &  Lorbeerblatt\\
    & Salz, Pfeffer\\
}

\preparation{%
    \step%
    Karotten, Sellerie, Pastinake, Zwiebeln und Lauch fein würfeln.
    
    \step%
    Olivenöl und Butter in einem Topf erhitzen und das Gemüse 6 Minuten anbraten, bis es leicht Farbe bekommt.
    
    \step%
    Währenddessen das Sojageschnetzeltes zubereiten, also in Gemüsebrühe und gegebenenfalls mit ein bisschen Sojasauce einlegen, ca 5min.

    \step%
    Sojageschnetzeltes klein schneiden und 4 Minuten mit dem Gemüse mitanbraten.

    \step%
    Tomaten blanchieren, also die Haut bisschen anscheiden und für 30 Sekunden in kochendes Wasser geben, danach in kaltem Wasser abschrecken. Die Haut kann dann leicht abgezogen werrden. Das Fruchtfleisch dann klein würfeln und zusammen mit dem Tomatenmark in den Topf geben.
    
    \step%
    Gewürze, sowie 3 EL Sojasahne hinzufügen und alles gut verrühren.
    
    \step%
    Auf niedriger Stufe 40 Minuten leicht köcheln lassen, bis die Sauce eindickt.
    
    \step%
    Zum Schluss 1 EL Sojasahne unterrühren und mit Salz und Pfeffer abschmecken.
}

\suggestion[Servieren]{%
    Mit 100 g trockener Pasta pro Portion servieren.
}

\suggestion[Tomaten]{%
    Die Tomaten sollten von guter Qualität sein, auch gut ist nicht nur eine Sorte, sondern verschiedene, kleinere Rispentomaten sind nach eigener Erfahrung etwas intensiver, die größeren etwas süßlicher.
}

\suggestion[Fleischersatz]{%
    Man könnte zwar auch Tofu oder fertiges Hack nehmen, jedoch finde ich das getrocknete Sojageschnetzeltes am besten, da dieses nach einer gewissen Kochzeit sehr fein wird und im Mund schön weich zergeht.
}

\hint{
    Das Gemüse kann man oft zusammen als feritgen Mix als Suppengemüse bekommen.
}

\end{recipe}